\documentclass{article}
\usepackage{amsmath, amssymb, amsthm}
\usepackage{hyperref}
\usepackage{geometry}
\geometry{margin=1in}

\title{Erd\H{o}s Problem \#11}
\date{Last updated: 2026-03-14}
\author{Source: erdosproblems.com}

\begin{document}
\maketitle

\noindent\textbf{Status:} OPEN \quad \textbf{No prize}

\noindent\textbf{Tags:} number theory, additive basis

\medskip

\noindent\textbf{Problem Statement:}

Is every large odd integer $n$ the sum of a squarefree number and a power of 2?




Odlyzko has checked this up to $10^7$. Hercher \cite{He24b} has verified this is true for all odd integers up to $2^{50}\approx 1.12\times 10^{15}$. Granville and Soundararajan \cite{GrSo98} have proved that this is very related to the problem of finding Wieferich primes , which are $p$ for which $2^{p-1}\equiv 1\pmod{p^2}$ - for example, if every odd integer is the sum of a squarefree number and a power of $2$ then a positive proportion of primes are non-Wieferich primes. Erd\H{o}s often asked this under the weaker assumption that $n$ is not divisible by $4$. Erd\H{o}s thought that proving this with two powers of 2 is perhaps easy, and could prove that it is true (with a single power of two) for almost all $n$. See also [9] , [10] , and [16] . This is mentioned in problem A19 of Guy's collection \cite{Gu04}.




References

[GrSo98] Granville, A. and Soundararajan, K., A Binary Additive Problem of Erd\H{o}s and the Order of $2$ mod $p^2$ . The Ramanujan Journal (1998), 283-298.

[Gu04] Guy, Richard K., Unsolved problems in number theory . (2004), xviii+437.

[He24b] C. Hercher, On the Sum of Squarefree Integers and a Power of Two . arXiv:2411.01964 (2024).

\noindent\textbf{Related OEIS sequences:} A001220, A377587


\bigskip
\noindent\small{Source: \url{https://www.erdosproblems.com/11}}
\end{document}
